\section*{Chapter \ref{ch:g3:files-and-folders} --- Files and Folders}

\begin{solution}{exo:g3:files-and-folders:1}
\code{map.png}: a picture. \code{notes.txt}: text. \code{dance.mp4}: a film.
\code{game.py}: a program written in Python. \code{song.mp3}: a sound.
\code{photo.jpg}: a picture. \code{list.txt}: text. \code{film.mp4}: a film.
\end{solution}

\begin{solution}{exo:g3:files-and-folders:2}
\code{Leo/Documents/School/Maths/sums.txt}.
\end{solution}

\begin{solution}{exo:g3:files-and-folders:3}
The \omterm{def:g3:files-and-folders:folder}{folder} \code{Pictures}, inside \code{Leo}; the \omterm{def:g3:files-and-folders:folder}{folder} \code{Documents},
inside \code{Leo}.
\end{solution}

\begin{solution}{exo:g3:files-and-folders:4}
(a)~Copy: she keeps her own \omterm{def:g3:files-and-folders:file}{file}. (b)~Move: the picture should be in one place
only, the right \omterm{def:g3:files-and-folders:folder}{folder}. (c)~Copy.
\end{solution}

\begin{solution}{exo:g3:files-and-folders:5}
Into the \omterm{def:g3:files-and-folders:bin}{bin}. It can be put back where it was, as long as the \omterm{def:g3:files-and-folders:bin}{bin} has not been
emptied.
\end{solution}

\begin{solution}{exo:g3:files-and-folders:6}
\begin{center}
\begin{forest}
  for tree={font=\footnotesize\sffamily, draw=black!60, rounded corners=1.5pt,
    fill=ocYellow!40, inner sep=2.5pt, l sep=4mm, s sep=2.5mm, edge={black!50}},
  where n children=0{fill=white, font=\footnotesize\ttfamily}{},
  [Mia [Drawings [cat.png] [house.png]] [Stories [dragon.txt]]]
\end{forest}
\end{center}
\end{solution}

\begin{solution}{exo:g3:files-and-folders:7}
Many plans are good. For example: a \omterm{def:g3:files-and-folders:folder}{folder} \code{School year} at the top, holding
one \omterm{def:g3:files-and-folders:folder}{folder} per subject (\code{Reading}, \code{Maths}, \code{Science},
\code{Art}\dots), and inside each subject one \omterm{def:g3:files-and-folders:folder}{folder} per project when there
are many \omterm{def:g3:files-and-folders:file}{files}. Each \omterm{def:g3:files-and-folders:file}{file} gets a name that says what it is.
\end{solution}

\begin{solution}{exo:g3:files-and-folders:8}
7 \omterm{def:g3:files-and-folders:folder}{folders} (\code{Leo}, \code{Documents}, \code{School}, \code{Poems},
\code{Maths}, \code{Pictures}, \code{Music}) and 8 \omterm{def:g3:files-and-folders:file}{files} (\code{poem.txt},
\code{rain.txt}, \code{sums.txt}, \code{letter.txt}, \code{cat.jpg},
\code{beach.png}, \code{rocket.png}, \code{song.mp3}).
\end{solution}

\begin{solution}{exo:g3:files-and-folders:9}
The \omterm{def:g1:what-is-a-computer:computer}{computer} thinks it is a picture and opens it with a picture viewer, which
cannot show it: it shows an error or nothing. The poem is not lost: the inside
of the \omterm{def:g3:files-and-folders:file}{file} has not changed. Renaming it back to \code{poem.txt} brings it back.
\end{solution}

\begin{solution}{exo:g3:files-and-folders:10}
Their paths are different, for example \code{Leo/Documents/notes.txt} and
\code{Leo/Music/notes.txt}, so the \omterm{def:g1:what-is-a-computer:computer}{computer} can tell them apart. In one \omterm{def:g3:files-and-folders:folder}{folder},
two \omterm{def:g3:files-and-folders:file}{files} with the same name would have the same path, and the \omterm{def:g1:what-is-a-computer:computer}{computer} would
not know which one you mean.
\end{solution}

\begin{solution}{exo:g3:files-and-folders:11}
Four moves: \code{Poems} into \code{Writing} (its poems go with it),
\code{story.txt} into \code{Writing}, \code{cat.jpg} and \code{dog.jpg} into
\code{Pictures}. Each of the four \omterm{def:g2:using-a-computer:desktop}{icons} that must change place needs one move,
so fewer is impossible.
\begin{center}
\begin{forest}
  for tree={font=\footnotesize\sffamily, draw=black!60, rounded corners=1.5pt,
    fill=ocYellow!40, inner sep=2.5pt, l sep=4mm, s sep=2.5mm, edge={black!50}},
  where n children=0{fill=white, font=\footnotesize\ttfamily}{},
  [Desktop, fill=black!8 [Pictures [cat.jpg] [dog.jpg]] [Writing [Poems, fill=ocYellow!40, font=\footnotesize\sffamily] [story.txt]]]
\end{forest}
\end{center}
\end{solution}
